% GENERATED FILE. Edit canonical lessons, then run npm run generate:books.
% canonical-source-hash: fnv1a64:87c895c6ae7864aa
% canonical-lessons: LA-C210-autem, LA-C210-vero, LA-C210-profecto, LA-C210-fere, LA-C210-tantum

\chapter{Moreover, Indeed, Surely, About, Only}
\label{ch:la-moreover-indeed-surely-about-only}
\hlchaptermodality{210}

\begin{chapteropening}
\textbf{By the end of this chapter:} \emph{I can say moreover, indeed, surely, about and only.}

Complete the last lesson of chapter 210: I can say moreover, indeed, surely, about and only.
\end{chapteropening}

\section[autem]{autem — moreover; however}

\label{lesson:LA-C210-autem}



\begin{quote}
Before the new one: say the Latin for in the way, then the Latin for all the way.
\end{quote}

\subsection*{You'll want to know: autem}
\textbf{autem} — \textquotedblleft{}moreover; however\textquotedblright{}.

\begin{cousinweb}
\textbf{autem} belongs with \emph{aut}, \textquotedblleft{}or\textquotedblright{} (Lewis \& Short), but how the two were formed is uncertain.
\end{cousinweb}

\begin{grammarlens}[title={What you've built}]
One more word for this chapter.
\end{grammarlens}

\subsection*{Your turn}
\begin{itemize}
\raggedright
  \item \emph{Say it:} \emph{autem}
  \item \emph{Say it:} \emph{autem}, once more
  \item \emph{Say it:} say \emph{autem}
  \item \emph{Recall:} say the Latin for in the way, then the Latin for all the way, then say \emph{autem} again
\end{itemize}

\begin{tcolorbox}[breakable,colback=teal!4,colframe=teal!35!black,title={Before you move on}]
What does autem mean? (\textquotedblleft{}Moreover; however\textquotedblright{}.) Say it once more.
\end{tcolorbox}
\section[vērō]{vērō — indeed; but}

\label{lesson:LA-C210-vero}



\begin{quote}
Before the new one: say the Latin for all the way, then the Latin for moreover; however.
\end{quote}

\subsection*{You'll want to know: vērō}
\textbf{vērō} — \textquotedblleft{}indeed; but\textquotedblright{}.

\begin{cousinweb}
\textbf{vērō} is an old case form of \emph{vērus}, \textquotedblleft{}true\textquotedblright{} (Lewis \& Short): in truth.
\end{cousinweb}

\begin{grammarlens}[title={What you've built}]
One more word for this chapter.
\end{grammarlens}

\subsection*{Your turn}
\begin{itemize}
\raggedright
  \item \emph{Say it:} \emph{vērō}
  \item \emph{Say it:} \emph{vērō}, once more
  \item \emph{Say it:} say \emph{vērō}
  \item \emph{Recall:} say the Latin for all the way, then the Latin for moreover; however, then say \emph{vērō} again
\end{itemize}

\begin{tcolorbox}[breakable,colback=teal!4,colframe=teal!35!black,title={Before you move on}]
What does vērō mean? (\textquotedblleft{}Indeed; but\textquotedblright{}.) Say it once more.
\end{tcolorbox}
\section[profectō]{profectō — surely}

\label{lesson:LA-C210-profecto}



\begin{quote}
Before the new one: say the Latin for moreover; however, then the Latin for indeed; but.
\end{quote}

\subsection*{You'll want to know: profectō}
\textbf{profectō} — \textquotedblleft{}surely\textquotedblright{}.

\begin{cousinweb}
\textbf{profectō} is \emph{prō factō}, \textquotedblleft{}for a fact\textquotedblright{} (Lewis \& Short): surely.
\end{cousinweb}

\begin{grammarlens}[title={What you've built}]
One more word for this chapter.
\end{grammarlens}

\subsection*{Your turn}
\begin{itemize}
\raggedright
  \item \emph{Say it:} \emph{profectō}
  \item \emph{Say it:} \emph{profectō}, once more
  \item \emph{Say it:} say \emph{profectō}
  \item \emph{Recall:} say the Latin for moreover; however, then the Latin for indeed; but, then say \emph{profectō} again
\end{itemize}

\begin{tcolorbox}[breakable,colback=teal!4,colframe=teal!35!black,title={Before you move on}]
What does profectō mean? (\textquotedblleft{}Surely\textquotedblright{}.) Say it once more.
\end{tcolorbox}
\section[ferē]{ferē — about, nearly}

\label{lesson:LA-C210-fere}



\begin{quote}
Before the new one: say the Latin for indeed; but, then the Latin for surely.
\end{quote}

\subsection*{You'll want to know: ferē}
\textbf{ferē} — \textquotedblleft{}about, nearly\textquotedblright{}.

\begin{cousinweb}
\textbf{ferē} has no sure origin; it is often linked through a root \emph{*dher-}, \textquotedblleft{}to hold fast\textquotedblright{} (EtymDB 2.1), with \emph{firmus}, \textquotedblleft{}firm\textquotedblright{}, but the link is uncertain.
\end{cousinweb}

\begin{grammarlens}[title={What you've built}]
One more word for this chapter.
\end{grammarlens}

\subsection*{Your turn}
\begin{itemize}
\raggedright
  \item \emph{Say it:} \emph{ferē}
  \item \emph{Say it:} \emph{ferē}, once more
  \item \emph{Say it:} say \emph{ferē}
  \item \emph{Recall:} say the Latin for indeed; but, then the Latin for surely, then say \emph{ferē} again
\end{itemize}

\begin{tcolorbox}[breakable,colback=teal!4,colframe=teal!35!black,title={Before you move on}]
What does ferē mean? (\textquotedblleft{}About, nearly\textquotedblright{}.) Say it once more.
\end{tcolorbox}
\section[tantum]{tantum — only; so much}

\label{lesson:LA-C210-tantum}



\begin{quote}
Before the new one: say the Latin for surely, then the Latin for about.
\end{quote}

\subsection*{You'll want to know: tantum}
\textbf{tantum} — \textquotedblleft{}only; so much\textquotedblright{}.

\begin{cousinweb}
\textbf{tantum} is the neuter of \emph{tantus}, \textquotedblleft{}so great\textquotedblright{}, used as an adverb (Lewis \& Short): so much, and only so much.
\end{cousinweb}

\begin{grammarlens}[title={What you've built}]
One more word for this chapter.
\end{grammarlens}

\subsection*{Your turn}
\begin{itemize}
\raggedright
  \item \emph{Say it:} \emph{tantum}
  \item \emph{Say it:} \emph{tantum}, once more
  \item \emph{Say it:} say \emph{tantum}
  \item \emph{Recall:} say the Latin for surely, then the Latin for about, then say \emph{tantum} again
\end{itemize}

\begin{tcolorbox}[breakable,colback=teal!4,colframe=teal!35!black,title={Before you move on}]
What does tantum mean? (\textquotedblleft{}Only; so much\textquotedblright{}.) Say it once more.
\end{tcolorbox}
